\item \model{MiniMax-M2}

\paragraph*{مقدمه و فلسفه بنیادین داده‌های پس‌آموزش عاملی}
سری مدل‌های \model{MiniMax-M2} (شامل نسخه‌های \en{M2}، \en{M2.5} و نگارش خودتکامل‌یافته \en{M2.7}) بر پایه دکترین \textbf{«فعال‌سازی‌های کمینه، آزادسازی بیشینه هوش در دنیای واقعی» (\en{Mini Activations Unleashing Max Real-World Intelligence})} پایه‌ریزی شده‌اند. معماری پایه مدل یک ترنسفورمر دیکودرمحور با ۲۲۹.۹ میلیارد پارامتر کل است که با بهره‌گیری از سازوکار خبرگان ریزدانه (\en{Fine-Grained MoE}) \cite{dai2024deepseekmoe}، در هر گام محاسباتی تنها \textbf{۹.۸ میلیارد پارامتر} را به ازای هر توکن فعال می‌نماید (فعال‌سازی ۸ خبره از میان ۲۵۶ خبره از طریق دروازه‌بانی سیگموئید با بایاس یادگیرنده بدون نیاز به زیان کمکی \cite{wang2024auxiliary}).

اگرچه مدل پایه روی ۲۹.۲ تریلیون توکن پیش‌آموزش دیده است، قابلیت‌های پیشرفته آن در سناریوهای استدلال بلندمدت، وظایف عاملی (\en{Agentic workflows}) و حل مسائل پیچیده صنعتی، تماماً در فاز پس‌آموزش (\en{Post-training}) شکل گرفته است. هسته اصلی مهندسی داده‌های پس‌آموزش این مدل، ارجحیت قاطع «کیفیت و تاییدپذیری پاداش مسیرها بر حجم داده‌ها» است. در این راستا، کلیه خطوط لوله بر پایه محیط‌های اجرایی واقعی و ایزوله (\en{Executable Workspaces}) و سیگنال‌های پاداش تاییدپذیر و منسجم با مصنوع نهایی (\en{Artifact-Aligned Verifiable Rewards}) سازمان‌دهی شده‌اند.

\paragraph*{خط لوله مهندسی نرم‌افزار سازمانی (\en{SWE-Scaling Pipeline})}
برای آموزش عامل‌های مهندسی نرم‌افزار سازمانی در مقیاس بالا و رفع نویز ذاتی مخازن گیت‌هاب، یک خط لوله شش‌مرحله‌ای خودکار و عامل‌محور توسعه یافته است \cite{jimenez2024swe}:
\begin{itemize}
    \item \textbf{گردآوری و پالایش درخواست‌های ادغام (\en{PR Collection \& Filtering}):} خزش در مخازن عمومی گیت‌هاب با مجوزهای آزاد، استخراج جفت‌های \en{PR} و \en{Issue} ادغام‌شده و فیلترگذاری مبتنی بر قواعد سخت‌گیرانه همراه با الزام حضور آزمون‌های واحد معتبر.
    \item \textbf{سنتز محیط‌های داکر چندزبانه توسط عامل (\en{Agent-Synthesized Docker}):} در زبان‌های کامپایلری مانند \en{Java}، \en{Go}، \en{Rust} و \en{C++}، یک عامل خودکار در یک چرخه تکرارشونده، اسکریپت‌های ساخت (\en{Build Scripts}) را تنظیم، در محیط کانتینری اجرا و خطاهای وابستگی را تا حصول محیط پایدار خود-اصلاح می‌کند. این فرآیند بیش از ۱۰ زبان برنامه‌نویسی را پوشش می‌دهد.
    \item \textbf{برچسب‌گذاری و مسیریابی وظایف (\en{PR Tagging \& Routing}):} تفکیک وظایف به دسته‌های رفع باگ (\en{Bug Fix})، افزودن قابلیت (\en{Feature Add})، بهینه‌سازی کارایی (\en{Performance Optimization}) و بازسازی کد به منظور طراحی توابع پاداش متمایز.
    \item \textbf{طراحی پاداش تاییدپذیر آزمون‌محور (\en{Verifiable Reward Construction}):} 
    در سناریوی رفع باگ، آزمون‌های شکست‌به-موفقیت (\en{Fail-to-Pass - F2P}) و موفقیت‌به-موفقیت (\en{Pass-to-Pass - P2P}) استخراج می‌شوند تا عدم بازگشت باگ‌های جانبی تضمین گردد. در افزودن قابلیت، آزمون‌های جدید برای ارزیابی مستقیم مؤلفه اضافه می‌شوند و در بهینه‌سازی کارایی، پایداری تفاوت کارایی قبل و بعد از اعمال تغییرات اندازه‌گیری می‌شود.
    \item \textbf{اعتبارسنجی وظیفه با مدل (\en{Model-based Task Validation}):} مدل ارزیاب، سازگاری میان شرح مسئله و آزمون‌ها را بررسی و اطلاعات ناقص را جهت خودبسنده‌سازی مسئله بازتولید می‌کند.
    \item \textbf{دگرگونی و افزایش تنوع دادگان (\en{Task Transformations}):} شامل تزریق عمدی باگ (\en{Bug Injection})، ادغام کامیت‌های مجاور جهت ساخت مسائل چندمرحله‌ای \cite{yang2024swesmith}، وارونه‌سازی مسئله به آزمون‌نویسی (\en{SWE-Test Conversion}) و وظایف مرور ایستا و تحلیل نقایص کد (\en{Code Review}).
\end{itemize}

\paragraph*{خط لوله توسعه نرم‌افزار با هدایت خبرگان (\en{AppDev Pipeline})}
توسعه اپلیکیشن از صفر مستلزم ارزیابی رفتار زمان اجرا و کیفیت طراحی است. در این حوزه از چارچوب ترکیبی زیر بهره گرفته شده است:
\begin{itemize}
    \item \textbf{سنتز کوئری با حضور خبره در حلقه (\en{Expert-in-the-Loop}):} مهندسان خبره قالب‌های متاروش (\en{Meta Queries}) شامل معماری‌ها و پشته‌های فناوری مدرن (نظیر \en{React + Zustand + Tailwind}) را تدوین می‌کنند. پس از نمونه‌برداری با دمای بالا، داده‌های تکراری با الگوریتم \en{MinHash} در زمان خطی $\mathcal{O}(N)$ حذف شده و توسط قاضی مدل زبانی بر اساس سه معیار عقلانیت پشته، امکان‌پذیری فنی و وضوح الزامات اعتبارسنجی می‌شوند.
    \item \textbf{تقطیر پرامپت (\en{Prompt Distillation Asymmetry}):} جهت حذف عادات نامطلوب طراحی (مانند پس‌زمینه‌های گرادیانی قالبی)، پرامپت‌های غنی به عامل آموزش می‌دهند که پیش از پیاده‌سازی، سند مشخصات بنویسد و فهرست \en{TODO} تهیه کند. در زمان نمونه‌برداری مسیرها، پرامپت کامل اعمال می‌شود اما در زمان آموزش، بخش‌هایی از آن تعمداً حذف می‌گردد تا مدل این رفتارهای بهینه را درونی‌سازی کند.
    \item \textbf{عامل به مثابه اعتبارسنج (\en{Agent-as-a-Verifier - AaaV}):} اپلیکیشن در محیط ایزوله مستقر شده و با کتابخانه \en{Playwright} در سه لایه سلسله‌مراتبی بررسی می‌شود:
    \begin{enumerate}
        \item \textit{لایه اجرا (\en{Execution Layer}):} بررسی نصب پکیج‌ها، بیلد بدون خطا و نبود خطاهای \en{JavaScript} در کنسول (عدم قبولی در این لایه موجب رد فوری است).
        \item \textit{لایه تعامل (\en{Interaction Layer}):} راستی‌آزمایی کارکرد دکمه‌ها، ارسال فرم‌ها و اجرای موفقیت‌آمیز سناریوهای کاربری.
        \item \textit{لایه زیبایی‌شناسی بصری (\en{Visual Aesthetics Layer}):} امتیازدهی به توازن چیدمان، هارمونی رنگ‌ها و ارگونومی رابط کاربری.
    \end{enumerate}
\end{itemize}

\paragraph*{چارچوب پایانه تعاملی (\en{Terminal-Gym})}
برای شبیه‌سازی سیستم‌عامل و تعاملات خط فرمان، پایگاه داده استک‌اورفلو پالایش شده و پست‌های اسکریپت‌پذیر، تاییدپذیر و سازگار با لینوکس استخراج شدند \cite{merrill2026terminal}. این خط لوله با تولید خودکار \en{Dockerfile}، آزمون‌های یکپارچه و انتزاع سرنخ‌های متنی، پلتفرم پایداری را برای سناریوهای پیشرفته نظیر \en{CVE-Factory} \cite{luo2026cve} و ابزارهای مهندسی خودکار فراهم می‌سازد.

\paragraph*{خط لوله‌های داده‌های همکاری عاملی (\en{Agentic Cowork})}
این بخش شامل چهار حوزه با محیط‌های کاری اجرایی و شبیه‌سازی‌شده است:
\begin{itemize}
    \item \textbf{جستجوی عمیق و وب‌پژوهی (\en{Deep Search}):} استفاده از راهبرد هدایت-و-بازنویسی (\en{Guide-and-Rewrite}) که در آن موجودیت‌های سوال پنهان می‌شوند تا مدل ناچار به ناوبری چندمرحله‌ای در وب گردد. مسیرها تنها در صورتی تایید می‌شوند که پاسخ نهایی مستقیماً بر مبنای شواهد بازیابی‌شده مستند باشد نه حافظه پارامتریک مدل.
    \item \textbf{وظایف اداری و کارکنان دانشی (\en{Office Tasks}):} با الهام از بنچمارک \en{GDPval} \cite{patwardhan2025gdpval}، وظایف متنوع حرفه‌ای سنتز شده و با فیلتر دقیق ارجاعات، از تولید داده‌های ساختگی و توهم جلوگیری می‌شود.
    \item \textbf{تحلیل مالی و صفحات گسترده (\en{Financial \& Spreadsheets}):} شامل دو شیوه سنتز داده است: سنتز مبتنی بر شواهد \cite{chen2026dive} که طی آن ابتدا ابزار مالی اجرا و سپس مسئله بر مبنای خروجی آن نگاشته می‌شود؛ و شیوه پیمایش کاربرگ \cite{liu2025webexplorer} که با اعمال تغییرات متوالی روی اکسل، وضعیت‌های میانی ذخیره و تطابق قطعی محاسبات ارزیابی می‌گردد.
    \item \textbf{تولید و ویرایش اسلاید (\en{Slide Generation}):} شامل طراحی افزایشی اسلایدها در سطوح عنصری و سندی با اعتبارسنجی بصری رندر توسط مدل‌های بینایی ماشین.
\end{itemize}

\paragraph*{داده‌های استدلال، مکالمه و نقش‌آفرینی}
برای وظایف استدلال عمیق ریاضی و علمی، مقیاس‌گذاری در سه محور دنبال شده است: تنوع مسائل (\en{Query-Side})، تولید چندین مسیر حل معتبر به ازای هر مسئله جهت تقویت تعمیم‌پذیری خارج از دامنه (\en{Response-Side}) و تنظیم پویا و بهینه نسبت ترکیب داده‌ها (\en{Training-Side}). در مکالمه همه‌منظوره، داده‌های زنجیره تفکر طولانی (\en{Long CoT}) با یک سیستم فایل تلفیق شده‌اند تا مهارت کار با اسناد تقویت شود. در بخش نقش‌آفرینی (\en{Role-Play}) \cite{minimax2026roleplay}، تولید شرطی روی فضای $\{\text{جهان‌ها}\} \times \{\text{داستان‌ها}\}$ تحت قید اولویت‌های کاربر مدل‌سازی شده و از الگوریتم‌های \en{RLHF} مبتنی بر استنتاج علّی و پایش آنتروپی برای رفع هک پاداش استفاده گردیده است.

\paragraph*{فرمول‌بندی تفکر متناوب در داده‌های \en{SFT}}
در فاز تنظیم دقیق با نظارت (\en{SFT})، داده‌ها به شیوه تفکر متناوب (\en{Interleaved Thinking}) مدل‌سازی می‌شوند که در آن توکن‌های استدلال $r_t$، اقدامات $a_t$ و مشاهدات محیطی $o_t$ در یک چرخه تکرار می‌شوند:
\begin{equation}
    \tau = (r_1, a_1, o_1, r_2, a_2, o_2, \dots, r_T, a_T, o_T)
\end{equation}
ویژگی کلیدی این داده‌ها، تداوم حالت استدلال (\en{Reasoning State Persistence}) است که در آن کل خروجی‌های تفکر نوبت پیشین در تاریخچه مکالمه $\mathcal{H}$ باقی می‌مانند:
\begin{equation}
    \mathcal{H}_{t+1} = \mathcal{H}_t \oplus [\text{assistant}(r_t, a_t)] \oplus [\text{tool}(o_t)]
\end{equation}
برخلاف ساختارهای فاقد حالت (\en{Stateless}) که تفکر پیشین حذف شده و مدل دچار خطای انباشته در استنتاج می‌گردد، این ساختار چرخه شناختی سه‌گانه «برنامه‌ریزی، اقدام و بازاندیشی» (\en{Plan-Act-Reflect}) را تقویت می‌نماید.

\paragraph*{مدل‌سازی داده‌ها در سیستم یادگیری تقویتی \en{Forge}}
تعامل عامل با محیط به عنوان یک فرآیند تصمیم‌گیری مارکوف $\mathcal{M} = (\mathcal{S}, \mathcal{A}, \mathcal{T}, \mathcal{R}, \gamma)$ با تابع انتقال $s_{t+1} = f_{\text{trans}}(s_t, a_t, o_t)$ مدل می‌شود. بهینه‌سازی از طریق الگوریتم \en{CISPO} \cite{minimax2025cispo} با نسبت اهمیت بریده‌شده نامتقارن صورت می‌پذیرد:
\begin{equation}
    J_{\text{CISPO}}(\theta) = \mathbb{E}_{\substack{(q,a) \sim \mathcal{D} \\ \{o_i\}_{i=1}^G \sim \pi_{\theta_{\text{old}}}}} \left[ \frac{1}{\sum_{i=1}^G |o_i|} \sum_{i=1}^G \sum_{t=1}^{|o_i|} \text{sg}\left(\hat{r}_{i,t}(\theta)\right) \hat{A}_{i,t} \log \pi_\theta(o_{i,t} \mid q, o_{i,<t}) \right]
\end{equation}
که در آن عملگر برش نسبت اهمیت به صورت زیر اعمال می‌شود:
\begin{equation}
    \hat{r}_{i,t}(\theta) = \text{clip}\left( \frac{\pi_\theta(o_{i,t} \mid q, o_{i,<t})}{\pi_{\theta_{\text{old}}}(o_{i,t} \mid q, o_{i,<t})}, \, 0, \, 1 + \epsilon_{\text{high}}^{\text{IS}} \right)
\end{equation}
تابع پاداش ترکیبی گام‌ها به فرم زیر تدوین شده است که پاداش ساختاری فرآیند، سرعت اجرا بر پایه تابع نزولی $h(\cdot)$ و صحت دستاورد نهایی را یکپارچه می‌سازد:
\begin{equation}
    r_t = \alpha \cdot r_t^{\text{process}} + \beta \cdot h\left(\frac{T_{\text{completion}}}{T_{\text{baseline}}}\right) + r_t^{\text{perf}}
\end{equation}
در زیرساخت توزیع داده‌ها، از استراتژی زمان‌بندی پنجره‌ای خروج به ترتیب ورود (\en{Windowed FIFO}) با اندازه پنجره $W = 0.3N$ و ادغام درخت پیشوند (\en{Prefix Tree Merging}) استفاده شده است. این نوآوری با محاسبه تنها یک‌باره بافت‌های طولانی مشترک در طول انتشار رو به جلو، ضمن انطباق ریاضی کامل، \textbf{شتاب‌بخشی تا ۴۰ برابر} را در پردازش داده‌های آموزش تقویتی به ارمغان می‌آورد.

\paragraph*{خودتکاملی داده‌ای و اسکلت‌بندی (\en{Self-Evolution})}
در مدل \model{MiniMax-M2.7}، چرخه تولید و تحلیل داده‌های شکست با اتکا بر سامانه تکرار مدل (\en{Model Iteration System}) خودکارسازی شده است. اسکلت تعاملی عامل کاملاً توسط نسخه درونی مدل بدون دخالت برنامه‌نویس انسانی پیاده‌سازی شده است. مدل به طور مداوم لاگ‌های آموزشی و ناهنجاری متریک‌ها را ریشه‌یابی کرده و با اصلاح کدهای پایپ‌لاین و فایل‌های تنظیمات، ۳۰ تا ۵۰ درصد بار کاری توسعه را خودکار نموده است. این سامانه در یک ارزیابی ۱۰۰ مرحله‌ای خودمختار، کارایی اسکلت ابزاری خود را تا ۳۰ درصد ارتقا داد.

\paragraph*{ارزیابی‌های آماری و تجربی فاز پس‌آموزش}
جدول \ref{tab:minimax_m2_sft_swa} نشان‌دهنده ارزیابی داده‌های \en{SFT} در مقیاس معماری مدل \en{M2} میان سازوکار توجه کامل (\en{Full Attention}) و پنجره لغزان ترکیبی (\en{Hybrid SWA}) است. نتایج اثبات می‌کند که با افزایش طول بافت به بیش از ۳۲ هزار توکن در وظایف عاملی، توجه کامل به دلیل پوشش جامع حافظه عملکرد بسیار برتری را به ثبت می‌رساند.

\begin{table}[htbp]
\centering
\caption{ارزیابی داده‌های \en{SFT} در مقیاس مدل \model{MiniMax-M2}: مقایسه توجه کامل و پنجره لغزان}
\label{tab:minimax_m2_sft_swa}
\resizebox{\textwidth}{!}{%
\begin{tabular}{llcc}
\toprule
\textbf{دامنه ارزیابی} & \textbf{بنچمارک (\en{Benchmark})} & \textbf{خط پایه (\en{Full Attention})} & \textbf{پنجره لغزان (\en{w/ SWA})} \\
\midrule
\multirow{5}{*}{\textbf{استدلال عمومی و دانش}} 
& \en{AIME 2025} & \textbf{86.7} & \textbf{86.7} \\
& \en{ARC-AGI-1} & 38.9 & \textbf{39.6} \\
& \en{GPQA-Diamond} & \textbf{75.3} & 72.7 \\
& \en{MMLU-Pro} & \textbf{80.5} & 80.1 \\
& \en{IFBench} & 23.1 & \textbf{27.2} \\
\midrule
\multirow{6}{*}{\textbf{وظایف عاملی بافت طولانی}} 
& \en{SWE-verified} & \textbf{54.7} & 50.2 \\
& \en{Terminal-Bench} & \textbf{26.7} & 23.8 \\
& \en{BrowseComp-zh} & \textbf{32.8} & 28.7 \\
& \en{GAIA-103} & \textbf{53.4} & 51.5 \\
& \en{XBench-ds} & 58.0 & \textbf{63.0} \\
& \en{$\tau^2$-Bench retail} & 62.3 & \textbf{67.5} \\
& \en{$\tau^2$-Bench telecom} & \textbf{32.5} & 21.0 \\
\bottomrule
\end{tabular}%
}
\end{table}

جدول \ref{tab:minimax_m27_frontier_eval} عملکرد نهایی مدل پس‌آموزش‌دیده \model{MiniMax-M2.7} (تنها با ۱۰ میلیارد پارامتر فعال در هر گام) را در مقایسه با قدرتمندترین مدل‌های مرزی جهان گزارش می‌کند. این داده‌ها حاکی از رقابت‌پذیری مستقیم این مدل با سامانه‌های بسیار حجیم‌تر تجاری در دامنه‌های کدنویسی، وب‌پژوهی، وظایف اداری و استدلال عمیق است.

\begin{table}[htbp]
\centering
\caption{مقایسه عملکرد \model{MiniMax-M2.7} در برابر مدل‌های مرزی متن‌بسته}
\label{tab:minimax_m27_frontier_eval}
\resizebox{\textwidth}{!}{%
\begin{tabular}{llcccccc}
\toprule
\textbf{دسته مهارتی} & \textbf{بنچمارک (\en{Benchmark})} & \textbf{\en{M2.7}} & \textbf{\en{M2.5}} & \textbf{\en{Claude Opus 4.6}} & \textbf{\en{Claude Sonnet 4.6}} & \textbf{\en{GPT 5.4}} & \textbf{\en{Gemini 3.1 Pro}} \\
\midrule
\multirow{6}{*}{\textbf{عاملی کدنویسی}} 
& \en{SWE-bench Pro} & 56.2 & 55.4 & 57.3 & 57.2 & \textbf{57.7} & 54.2 \\
& \en{SWE-bench Multilingual} & 76.5 & 74.1 & \textbf{77.8} & 75.9 & 70.5 & – \\
& \en{Multi-SWE-bench} & \textbf{52.7} & 51.3 & 50.3 & 51.0 & 49.0 & – \\
& \en{NL2Repo} & 39.8 & 26.6 & 43.7 & 43.3 & \textbf{46.8} & 35.9 \\
& \en{Terminal-Bench 2.0} & 57.0 & 51.7 & 65.4 & 59.1 & \textbf{75.1} & 68.5 \\
& \en{MLE Bench Lite} & 66.6 & 51.5 & \textbf{75.7} & 72.7 & 71.2 & 66.6 \\
\midrule
\multirow{2}{*}{\textbf{توسعه کاربردی}} 
& \en{VIBE-Pro} & 55.6 & 54.2 & 55.6 & \textbf{56.1} & – & 41.0 \\
& \en{HyperTask} & 67.6 & 59.4 & \textbf{75.7} & 74.1 & – & 50.9 \\
\midrule
\multirow{3}{*}{\textbf{جستجوی عمیق}} 
& \en{BrowseComp} & 77.8 & 76.3 & 84.0 & 74.7 & 82.7 & \textbf{85.9} \\
& \en{Wide Search} & 75.2 & 70.3 & \textbf{79.4} & 75.8 & 77.9 & – \\
& \en{RISE} & 64.3 & 50.2 & \textbf{68.5} & 58.8 & 63.3 & – \\
\midrule
\multirow{5}{*}{\textbf{ابزار و امور اداری}} 
& \en{GDPval-AA} & 50.0 & 35.0 & 55.0 & 57.0 & \textbf{58.0} & 41.0 \\
& \en{Toolathlon} & 46.3 & 38.3 & 47.2 & 44.8 & \textbf{54.6} & 48.8 \\
& \en{MM Claw} & 62.7 & 57.6 & \textbf{75.4} & 64.2 & 73.6 & 61.8 \\
& \en{MEWC v2} & 63.3 & 49.8 & 62.0 & \textbf{77.2} & 76.5 & 42.2 \\
& \en{Finance Modeling Pro} & 57.0 & 33.8 & 69.0 & 66.2 & \textbf{75.3} & 35.6 \\
\midrule
\multirow{7}{*}{\textbf{استدلال و دانش}} 
& \en{AIME 2026} & 94.2 & 87.2 & 92.5 & 92.7 & \textbf{97.0} & 88.7 \\
& \en{GPQA-Diamond} & 89.8 & 85.2 & 89.6 & 87.5 & 92.0 & \textbf{94.1} \\
& \en{SciCode} & 47.0 & 43.0 & 51.9 & 46.8 & 56.6 & \textbf{58.9} \\
& \en{IFBench} & 76.0 & 72.0 & 53.1 & 56.6 & 73.9 & \textbf{77.1} \\
& \en{AA-LCR} & 72.0 & 65.0 & 70.7 & 70.7 & \textbf{74.0} & 72.7 \\
& \en{HLE} & 28.0 & 19.0 & 36.7 & 30.0 & 41.6 & \textbf{44.7} \\
& \en{MMLU-Pro} & 81.8 & 85.2 & 89.1 & 87.3 & 87.5 & \textbf{91.2} \\
\bottomrule
\end{tabular}%
}
\end{table}

بررسی مسیر تکامل درون‌نسلی مدل نشان می‌دهد که ورود داده‌های تخصصی جدید در نگارش‌های \en{M2.5} و \en{M2.7} بیشترین جهش‌های خطی را ایجاد کرده است؛ من‌جمله ارتقای چشمگیر در بنچمارک \en{BrowseComp} از $44.0$ در مدل اولیه \en{M2} به $76.3$ در \en{M2.5} و سپس $77.8$ در \en{M2.7} ($+33.8$ امتیاز رشد)، جهش نرخ مدال در \en{MLE Bench Lite} از $40.0$ به $66.6$ ($+26.6$ امتیاز رشد که در آن بهترین اجرای مدل موفق به کسب ۹ مدال طلا، ۵ نقره و ۱ برنز شد) و رشد امتیاز وظایف اسناد اداری در \en{GDPval-AA} از $34.0$ به $50.0$ ($+16.0$ امتیاز رشد).